\section{Limitations}
\label{sec:limitations}

\subsection*{II-A limitations (preserved, amended for Oracle semantics)}
\begin{itemize}
 \item \textbf{Single validated real workload/trace for II-A.} The II-A
   replay, comparator-gap, sensitivity, and decomposition analyses are
   computed against one captured coding trace (one workload, one
   30B-A3B model instance, one capture budget). Replay fidelity
   (Section~\ref{sec:replay-validation}) is demonstrated for this trace
   only; it is not evidence of fidelity for other workloads, models, or
   capture configurations.
 \item \textbf{Offline replay, not live measurement.} II-A comparisons are
   offline replays against a fixed, already-captured trace. They do not
   measure live runtime behavior under the compared policies themselves.
 \item \textbf{No latency or throughput claim.} Miss bytes are a proxy for
   expert data that would need to be serviced under the replay model. No
   latency, throughput, tokens/second, or SSD-traffic measurement is made
   or implied by II-A or II-B.
 \item \textbf{No causal predictor exists yet.} The Belady comparator gap
   (Section~\ref{sec:recoverable-envelope}) is a measured separation from
   $O_B$, not a demonstrated causal recovery and not a proven bound on
   causal recovery. II-B characterizes routing structure but does not
   evaluate a causal predictor or recover any fraction of that gap.
 \item \textbf{Comparator-objective limitation (\textsc{AUDIT-ORACLE-001},
   CASE~C).} Expert cache objects are $(\ell,e)$ pairs with two byte-size
   classes across layers in the tested trace. The implemented
   \textsc{Oracle-Global-Belady} comparator ($O_B$) uses furthest next use
   only; object byte size does not enter victim ranking. Miss-count and
   miss-byte optimality are not proven under the actual variable-size,
   byte-capacity, multi-key replay semantics. Consequently $B-O_B$ is not
   interpreted as a proven upper or lower bound on causal recovery
   $B-C^*$; $O^* \le O_B$ is known, but $O_B$ and $C^*$ have no guaranteed
   ordering.
 \item \textbf{Statistical, not exact, shard mapping.} The
   Runtime-Sharded-LRU baseline matches the real runtime's \emph{shard
   count} (16 shards at $4$~GiB) but not its exact address-derived
   shard-to-key mapping, which is not reproducible offline.
 \item \textbf{Finite, predeclared budget grid.} The II-A budget sweep
   tests five fixed points ($1.5, 2, 4, 8, 16$~GiB). Behavior between
   tested points and the exact mathematically minimal policy-convergence
   budget are not characterized.
 \item \textbf{LFU auxiliary sweep scope.} The LFU-vs-Global-LRU ranking
   across the R2-03 budget grid is canonical auxiliary R2-03 evidence from
   the same single canonical DEV trace and tested budget grid. It is not
   held-out evaluation and does not establish causal-baseline ordering
   across other workloads, models, or capacities. R2-09A remains necessary.
 \item \textbf{No shard-constrained clairvoyant comparator.} The II-A
   structural and comparator decomposition uses Global-LRU as a shared
   pivot rather than a true sharded-clairvoyant comparison.
\end{itemize}

\subsection*{II-B limitations}
\begin{itemize}
 \item \textbf{Model, quantization, and runtime scope.} II-B evidence is
   limited to Qwen3-30B-A3B-Q4\_K\_M under one consumer-machine runtime
   protocol and one expert-budget configuration context for the tested
   captures. No cross-model, cross-quantization, cross-machine, or
   cross-runtime generalization is measured.
 \item \textbf{Workload scope.} R2-08 tests four workload families
   (coding, reasoning, summarization, instruction), two replicates each.
   The \texttt{WORKLOAD\_DEPENDENCE\_DOMINANT} classification is
   conditional on this tested workload set and protocol.
 \item \textbf{Deterministic decoding protocol.} R2-08 uses
   temperature~$=0.0$, top-$k=1$. Same-workload replicate JSD $=0.0$ is
   therefore interpreted only as reproducibility at the reported
   resolution under deterministic decoding, not as general evidence of
   zero stochastic run-to-run variability.
 \item \textbf{Finite matched horizon.} R2-08 comparisons use matched
   decode horizon $H=388$ with 8 temporal windows of length 48. Longer
   horizons and other windowing choices are not evaluated.
 \item \textbf{Layer-local descriptive analysis.} II-B metrics are
   layer-local and descriptive. No cross-layer primary transition claim is
   made from R2-06, and no causal or mechanistic explanation of routing is
   established.
 \item \textbf{No predictive or runtime evaluation.} II-B includes no
   causal predictor, no predictor accuracy/precision/recall/F1 measurement,
   no prefetch or admission mechanism, no retention-policy replacement, no
   runtime speedup or latency/throughput measurement, and no
   comparator-gap recovery measurement. The comparator-relative metric in
   Eq.~\eqref{eq:iib-iic-recovery} remains \notyetmeasured{}.
 \item \textbf{R2-08 capture-provenance representation.} New R2-08-managed
   capture cells record \texttt{capture\_run\_id}/\texttt{capture\_source}
   as null while preserving trace path, trace SHA, configuration, model
   and runtime SHAs, validation state, and canonical R2-08 run identity.
   This is documented as an intentional representation for internal R2-08
   captures; no capture IDs are invented and no scientific traces are
   modified.
\end{itemize}

\subsection*{DEV/EVAL, EXP-025, and II-C evaluation transparency}
\begin{itemize}
 \item \textbf{DEV characterization evidence.} The canonical R2-01 coding
   trace is documented as a \texttt{dev} split in inspected provenance, and
   the inspected R2-08 workload captures are also documented as
   \texttt{dev}-split characterization captures. Research~II-A and II-B are
   characterization/development tranches, not held-out predictive
   evaluation. Research~II-C's R2-09B contract subsequently froze DEV/EVAL
   roles before any EVAL capture, and R2-10/R2-BRIDGE-01/R2-11/R2-12 each
   captured and sealed their held-out instances before scoring; see
   Section~\ref{sec:iic} for results and Section~\ref{sec:iic-limitations}
   below for this stage's own limitations.
 \item \textbf{EXP-025 applicability.} \textsc{AUDIT-ORACLE-001} classifies
   EXP-025 as a stream/distribution-dependent performance model (P2) whose
   transferability to arbitrary policy-induced miss streams is
   \textsc{UNSUPPORTED} without separate validation. Comparator miss-byte
   reductions must not be converted directly into latency or avoided
   seconds. Future modeled comparisons, if used, must follow
   \begin{equation*}
   \Delta T_{\mathrm{model}} = \mathrm{PerformanceModel}(\mathrm{MissStream}_{B_\mathrm{best}})
   - \mathrm{PerformanceModel}(\mathrm{MissStream}_P)
   \end{equation*}
   with justified scope, and must not be equated with measured end-to-end
   latency, time-to-first answer, or throughput.
 \item \textbf{No numeric practical-effect threshold.} No numeric
   practical-effect threshold is preregistered by this amendment. Any
   future threshold must be derived and frozen in R2-09B using a metric
   justified by the evaluation/performance-model contract.
\end{itemize}

\subsection*{II-C limitations}
\label{sec:iic-limitations}
\begin{itemize}
 \item \textbf{One tested predictive-retention realization.} Only one
   causal predictor (\texttt{decayed\_transition}, decay $0.99$) driving
   one reactive retention-policy realization (\texttt{PredictedProtectPolicy})
   was built and evaluated. Negative results at R2-12 are results about
   this specific realization, not a general statement about predictive
   retention as a design category.
 \item \textbf{Hardware/model/workload scope.} All of R2-09 through R2-12
   were evaluated on one model family (Qwen3-30B-A3B), one consumer host,
   and workload instances drawn from four families (coding, instruction,
   reasoning, summarization). No cross-model, cross-quantization,
   cross-machine, or cross-runtime generalization is claimed.
 \item \textbf{Replay-level benefit does not imply runtime benefit.} R2-11's
   positive replay-level result (\textsc{r2\_11\_pass}) and R2-12's negative
   real-runtime result (\textsc{R2\_12\_FAIL}) are not in tension: they
   measure different quantities (miss bytes per decode step vs. decode
   tokens per second) under different experimental apparatus (offline
   replay vs. the real runtime). Neither bounds the other.
 \item \textbf{Runtime failure does not invalidate predictor quality.} R2-10's
   held-out-confirmed prediction-quality margin (Section~\ref{sec:r2-10})
   is not weakened, retracted, or reclassified by R2-12's negative result.
 \item \textbf{Secondary budgets are not primary decision endpoints.} The
   1.5~GiB and 8~GiB R2-12 budgets are reported descriptively only; neither
   can, under the frozen decision rule, override the primary 4~GiB
   classification, regardless of their own effect sizes or direction.
 \item \textbf{Telemetry gap limits causal attribution at 8~GiB.} The real
   8~GiB canonical records carry no retention-traffic telemetry
   (\texttt{expert\_bytes\_read}/\texttt{expert\_read\_count}/
   \texttt{expert\_miss\_bytes} and every \texttt{expert\_io\_best\_effort}
   sub-field are \texttt{None} for both policies). Any non-retention-overhead
   explanation of the 8~GiB delta is therefore reported as an unverified
   hypothesis, never as a finding.
 \item \textbf{Cross-budget ordering is descriptive.} The ordering of
   measured deltas across 1.5/4/8~GiB (smallest to largest in magnitude) is
   reported as an observation over three tested points, not as evidence of
   a monotonic causal relationship between memory budget and policy
   overhead.
 \item \textbf{Results do not establish universal impossibility.} A negative
   result for one tested reactive-retention realization does not establish
   that no predictive mechanism (e.g., active prefetch, scheduling, or
   compute/I/O overlap -- Section~\ref{sec:iic-related-work}) can convert
   predictive information into real-runtime throughput gain in this problem
   class.
 \item \textbf{Host/environment constraints are exactly those frozen in the
   protocol, no more.} R2-12's machine-state and readiness gates
   (AC power, operator-attested clean environment, thermal cooldown policy,
   power mode) are exactly the conditions frozen in
   \texttt{machine\_state\_contract.json} and \texttt{execution\_plan.json}
   before acquisition. No condition not formally frozen there (e.g., a
   specific battery-percentage or ambient-temperature threshold used
   elsewhere in the Argine programme) is retroactively applied as an
   exclusion criterion to R2-12's evidence.
 \item \textbf{Runtime conformance was reached iteratively.} The runtime
   binary realizing the retention-policy integration for R2-12 required
   several realization/conformance iterations before reaching an accepted
   conformance state; it is not byte-identical to the R2-01 validated
   replay-comparison runtime. This is disclosed rather than assumed from a
   single attempt.
\end{itemize}
